% Chapter 5: Experiments, Results and Discussion
\chapter{Experiments, Results and Discussion}\label{ch:experiments}

\guide{Describe the setup, present the results objectively, compare them with
existing methods, and then interpret them. The discussion should say what the
results mean for each research question.}

This chapter reports the experiments in the order in which they were run,
because later experiments were planned on the basis of earlier outcomes. After
the setup (\cref{sec:setup}) and the validation of the evaluation code
(\cref{sec:metric-validation}), the two preliminary gates are reported: the
choice of the video encoder (\cref{sec:g1}) and the test of predictive surprise
as a fixation cue (\cref{sec:g2}). The main results on DHF1K, the transfer to
eight further test sets, qualitative examples and the computational cost follow
in \cref{sec:results}, and the ablations in \cref{sec:ablation}. The chapter
ends with a discussion organised by research question (\cref{sec:discussion})
and an assessment of the threats to validity (\cref{sec:threats}). Unless stated
otherwise, every number in this chapter was produced by the code and evaluator
described in \cref{ch:methodology}, and every test set was scored once per
reported model.

%----------------------------------------------------------------------------
\section{Experimental Setup}\label{sec:setup}

All experiments ran on a single workstation (\cref{tab:setup}), one job at a
time. A monitoring process checked the GPU temperature, hardware events, disk
space and memory every minute and could stop training cleanly; it was never
triggered.

\begin{table}[htbp]
  \centering
  \caption{Experimental environment}\label{tab:setup}
  \begin{tabular}{@{}ll@{}}
    \toprule
    \textbf{Item} & \textbf{Specification} \\
    \midrule
    \ac{GPU}            & NVIDIA RTX A5000, 24\,GB \\
    CPU / RAM           & Intel Xeon w7-2475X (20 cores) / 256\,GB \\
    Operating system    & Windows 11 Pro \\
    Language            & Python 3.12 \\
    Framework           & PyTorch 2.14 with CUDA 13.0~\cite{paszke2019pytorch} \\
    Model loading       & Transformers 5.22~\cite{wolf2020transformers} \\
    Other libraries     & NumPy~\cite{harris2020array}, OpenCV, SciPy \\
    \bottomrule
  \end{tabular}
\end{table}

\subsection{Implementation Details}

The settings of \cref{tab:hyper} were used throughout. Training one
three-epoch DHF1K model took about 15 minutes of training time with the
ViT-L encoder (4.7 minutes per epoch of 4,320 clips) and about 31 minutes with
ViT-g (9.9 minutes per epoch), plus the holdout evaluation after each epoch.
ViT-g fitted in 10\,GB of GPU memory at a batch size of eight. In total, the
experiments of this chapter comprise 13 DHF1K training runs for the gates, the
main program and its ablations, 7 further DHF1K runs for the ViT-g, 600-video
and 32-frame experiments, and 125 fine-tuning runs on the other datasets
(including the 16 runs with audio). Holdout lists, configuration files and
the log of every decision are part of the code accompanying the thesis.

%----------------------------------------------------------------------------
\section{Validation of the Evaluator}\label{sec:metric-validation}

Because the metrics were implemented from scratch, the evaluator was
validated before any SURGE-D model was scored (\cref{tab:rules}). The
reference maps were the DHF1K validation predictions of
SalFoM~\cite{moradi2024salfom}, produced on the same machine with the authors'
released weights, for which per-frame scores computed with the evaluation code
released by the TMFI-Net authors~\cite{zhou2023transformer} were available.

The evaluator reproduced the reference scores: CC 0.5639 against 0.5642
($\Delta=-0.0003$) and NSS 3.3264 against 3.3268 ($\Delta=-0.0004$), well within
the tolerance of the rule. The SIM, KLD and AUC-J values of the reference code
use other normalisations and could not be compared directly; against a second,
independently written evaluator that follows the MIT benchmark conventions,
CC, NSS, SIM and KLD agreed per video to within $5\times10^{-5}$ and AUC-J agreed
in the mean. Two deviations from the reference procedures were found and
corrected during this check, before any model was scored: the first AUC-J
implementation lacked the tie-breaking jitter of the reference code (which
lowered AUC-J by 0.010 on quantised maps), and the first shuffled AUC pooled all
negatives instead of drawing them per split. With the corrected evaluator,
SalFoM's validation maps score CC 0.5639, NSS 3.326, SIM 0.4130, KLD 1.297,
AUC-J 0.9243 and sAUC 0.7406. These are the SalFoM numbers used in all
comparisons below. They are slightly lower than SalFoM's published scores
because they are computed on the validation set rather than on the official
test set, on which the authors report CC 0.569 and NSS
3.353~\cite{moradi2024salfom}.

%----------------------------------------------------------------------------
\section{Selection of the Video Encoder (Gate G1)}\label{sec:g1}

The proposal had chosen V-JEPA~2.1~\cite{murlabadia2026vjepa21}, which extends
the V-JEPA~2 objective with a loss on the visible tokens and with
self-supervision at intermediate layers, specifically to obtain denser
features. Its ViT-L model is distilled from a two-billion-parameter teacher and
was cooled down at $384\times384$ pixels~\cite{murlabadia2026vjepa21}\footnote{Checkpoint
\texttt{vjepa2\_1\_vitl\_dist\_vitG\_384} from the official repository.}.
Gate G1 compared three ViT-L encoders under identical conditions: V-JEPA~2,
V-JEPA~2.1 and UMT-L~\cite{li2023unmasked}, the encoder family used by SalFoM.
Each was frozen and trained with the same decoder for ten epochs with seed 0,
all V-JEPA models at $256\times256$ pixels. UMT-L received $224\times224$ input
because its position table is fixed to that size and to 16 frames.

\begin{table}[htbp]
  \centering
  \caption{Gate G1: frozen ViT-L encoders with the same decoder, internal DHF1K
  holdout (60 videos, every fifth frame, no blur), best of ten epochs, seed 0}\label{tab:g1}
  \begin{tabular}{@{}lccccc@{}}
    \toprule
    \textbf{Encoder} & \textbf{Best epoch} & \textbf{CC}$\uparrow$ & \textbf{NSS}$\uparrow$
      & \textbf{SIM}$\uparrow$ & \textbf{KLD}$\downarrow$ \\
    \midrule
    V-JEPA~2 ViT-L ($256^2$)   & 3 & \textbf{0.5782} & \textbf{3.314} & \textbf{0.4432} & \textbf{1.218} \\
    V-JEPA~2.1 ViT-L ($256^2$) & 2 & 0.5692 & 3.295 & 0.4377 & 1.245 \\
    UMT-L ($224^2$)~\cite{li2023unmasked} & 2 & 0.5464 & 3.101 & 0.4157 & 1.323 \\
    \bottomrule
  \end{tabular}
\end{table}

\Cref{tab:g1} shows that V-JEPA~2 led on all four metrics. V-JEPA~2.1 trailed by
0.009 CC, outside the margin of 0.005 within which the proposal's choice would
have been kept, so V-JEPA~2 was adopted for all later experiments. UMT-L trailed
by 0.032 CC. With one seed per encoder, part of the gap between the two
V-JEPA variants may be seed noise; the gap to UMT-L is several times larger
than the seed-to-seed spread observed later (0.003 CC). The ranking of the two
V-JEPA variants runs against the expectation that denser features help a dense
task. Two explanations are compatible with the data and were not separated:
the decoder reads four intermediate layers, which may already supply the local
detail that the last layer of V-JEPA~2 lacks, and the V-JEPA~2.1 model was used
at a lower resolution than the one at which it was cooled down.

The learning curves in \cref{fig:g1} were the more consequential outcome of this
gate. For every encoder, holdout CC peaked in the second or third epoch and
then declined steadily, while training CC kept rising. A frozen encoder with a
7-million-parameter decoder thus overfits the 540 training videos after a few
passes. On the basis of these curves, and before any further run, the schedule
of the main program was shortened to three epochs (\cref{sec:training}).

\begin{figure}[htbp]
  \centering
  \includegraphics[width=\textwidth]{g1_curves.pdf}
  \caption{Learning curves of gate G1. (a) Holdout CC after each epoch; circles
  mark the best epoch. (b) Mean CC on the training clips of each epoch. Holdout
  accuracy peaks early and then declines while training accuracy keeps rising,
  for all three encoders}\label{fig:g1}
\end{figure}

%----------------------------------------------------------------------------
\section{Predictive Surprise as a Fixation Cue (Gate G2)}\label{sec:g2}

Gate G2 tested the hypothesis that motivated the original proposal: that the
error of V-JEPA~2's pretrained predictor indicates where people look
(\cref{sec:surprise}). No training is involved; the surprise maps and the
reference maps were scored directly against the fixations of the 60 holdout
videos.

\begin{table}[htbp]
  \centering
  \caption{Gate G2: zero-shot fixation prediction by surprise and reference maps
  on the DHF1K holdout (60 videos), at about 4 and 15 sampled frames per
  second}\label{tab:g2}
  \begin{tabular}{@{}lcccc@{}}
    \toprule
    & \multicolumn{2}{c}{\textbf{NSS}$\uparrow$} & \multicolumn{2}{c}{\textbf{CC}$\uparrow$} \\
    \cmidrule(lr){2-3}\cmidrule(l){4-5}
    \textbf{Map} & 4\,fps & 15\,fps & 4\,fps & 15\,fps \\
    \midrule
    Raw surprise, \cref{eq:surprise}        & $-0.190$ & $-0.502$ & $-0.043$ & $-0.099$ \\
    Contextual surprise (true $-$ other past) & $0.116$ & $-0.366$ & $0.018$ & $-0.072$ \\
    Token atypicality                        & $0.454$ & $0.650$ & $0.081$ & $0.123$ \\
    Frame difference                         & $0.810$ & $0.803$ & $0.149$ & $0.148$ \\
    Empirical centre prior                   & $1.464$ & $1.479$ & $0.289$ & $0.292$ \\
    \bottomrule
  \end{tabular}
\end{table}

The outcome was unambiguous (\cref{tab:g2}). Neither surprise variant beat token
atypicality or frame difference at either frame rate, and the raw surprise was
\emph{anti}-correlated with fixations: its NSS is negative, so observers looked
on average at regions that the predictor found \emph{easier} than average to
predict. Information gain over the centre prior was negative for every map.
The gate therefore failed, and the surprise module was never trained.

A preliminary diagnostic, run before the gate on eleven clips without gaze data
(ten Kinetics validation clips~\cite{carreira2017quo} and one animated film
trailer), helps to explain this result. When the true past of a clip was
replaced by another clip's past, the mean prediction error rose in every clip
(from 0.644 to 0.731 at about 4\,fps), so the predictor does use the context.
Yet the spatial pattern of the error correlated at $r\approx0.60$ with token
atypicality and at $r\approx0.62$ with the error map computed from the
\emph{wrong} past, and it correlated negatively with a central Gaussian
($r\approx-0.11$). The error map therefore mainly reflects how unusual each
target region is, largely independently of the past, and it is lower at the
centre of the frame, where observers tend to look. The released training
configuration of V-JEPA~2 offers a plausible reason: its masks span the whole
clip in time\footnote{\texttt{temporal\_scale: [1.0, 1.0]} in the released
pre-training configuration files of the official repository.}, so the predictor
was trained to fill in masked scene content rather than to extrapolate into the
future. Its error is consequently not an expectation-violation signal in the
sense of Bayesian surprise~\cite{itti2009bayesian}.

%----------------------------------------------------------------------------
\section{Results}\label{sec:results}

\subsection{Quantitative Results on DHF1K}\label{sec:main-results}

The main program trained the frozen-encoder model (labelled B1 in the research
log) with three seeds. The LoRA variant lost to it in all three seeds
(\cref{sec:ablation}), so the frozen model became the main model. After the
ViT-g encoder passed its own gate, the same recipe was repeated with ViT-g.
\Cref{tab:dhf1k} reports all four configurations on the DHF1K validation set,
next to SalFoM scored by the same evaluator and to published validation results
of other methods.

\begin{table}[htbp]
  \centering
  \caption{Results on the DHF1K validation set (100 videos, 60,273 frames). The
  upper block was scored by the evaluator of this thesis; the lower block is
  copied from the literature and was computed with other evaluation code. Best
  value per column in bold}\label{tab:dhf1k}
  \small
  \begin{tabular}{@{}lcccccc@{}}
    \toprule
    \textbf{Model} & \textbf{CC}$\uparrow$ & \textbf{NSS}$\uparrow$ & \textbf{SIM}$\uparrow$
      & \textbf{KLD}$\downarrow$ & \textbf{AUC-J}$\uparrow$ & \textbf{sAUC}$\uparrow$ \\
    \midrule
    SalFoM~\cite{moradi2024salfom} (released weights) & 0.5639 & 3.326 & 0.4130 & 1.297 & 0.9243 & 0.7406 \\
    SURGE-D ViT-L, single ($\sigma=12$)    & 0.5836 & 3.349 & 0.4492 & 1.222 & 0.9283 & 0.7636 \\
    SURGE-D ViT-L, 3-seed ensemble ($\sigma=8$) & 0.5910 & 3.418 & 0.4486 & 1.187 & 0.9302 & 0.7614 \\
    SURGE-D ViT-g, single ($\sigma=8$)     & 0.6003 & 3.457 & 0.4531 & 1.170 & 0.9308 & \textbf{0.7709} \\
    SURGE-D ViT-g, 3-seed ensemble ($\sigma=8$) & \textbf{0.6029} & \textbf{3.486} & \textbf{0.4574}
      & \textbf{1.160} & \textbf{0.9314} & 0.7680 \\
    \midrule
    \multicolumn{7}{@{}l}{\emph{Published, other evaluators (as reported in~\cite{girmaji2025minimalistic})}} \\
    TMFI-Net~\cite{zhou2023transformer}      & 0.554 & 3.201 & 0.428 & -- & 0.924 & -- \\
    THTD-Net~\cite{moradi2024transformer}    & 0.553 & 3.188 & 0.425 & -- & 0.924 & -- \\
    ViNet-E~\cite{girmaji2025minimalistic}   & 0.549 & 3.134 & 0.409 & -- & 0.922 & -- \\
    STSANet~\cite{wang2023spatio}            & 0.539 & 3.082 & 0.411 & -- & 0.920 & -- \\
    DiffSal~\cite{xiong2024diffsal}          & 0.533 & 3.066 & 0.405 & -- & 0.918 & -- \\
    \bottomrule
  \end{tabular}
\end{table}

\paragraph{Comparison with SalFoM.} SalFoM is the closest prior work: it was the
first video saliency model built on a video foundation encoder, it reported the
best DHF1K test results among the methods it compared
with~\cite{moradi2024salfom}, and its released weights allowed a comparison
under an identical protocol. \Cref{tab:dhf1k-boot} gives the paired bootstrap
intervals. The ViT-g model is better than SalFoM on all six metrics as a single
model and as an ensemble, with every interval excluding zero; the CC gain of
the single model is $+0.036$ $[+0.026, +0.048]$. The ViT-L ensemble is also
better on all six metrics, and the single ViT-L model on five, its NSS
difference not being significant. The per-video differences in \cref{fig:valdelta} show that
the gain is broad rather than driven by a few videos: the ViT-g ensemble has
the higher CC on 79 of the 100 videos.

\begin{table}[htbp]
  \centering
  \caption{Paired bootstrap differences on DHF1K validation (100 videos, 10,000
  resamples), with 95\,\% intervals; differences whose interval excludes zero are
  in bold. The intervals for all six metrics are given in
  \cref{tab:app-boot}}\label{tab:dhf1k-boot}
  \small
  \begin{tabular}{@{}lcc@{}}
    \toprule
    \textbf{Comparison} & $\Delta$\textbf{CC} [95\,\% CI] & $\Delta$\textbf{NSS} [95\,\% CI] \\
    \midrule
    ViT-L single $-$ SalFoM           & $\mathbf{+0.020}$ [+0.010, +0.031] & +0.023 [$-$0.050, +0.097] \\
    ViT-L ensemble $-$ SalFoM         & $\mathbf{+0.027}$ [+0.017, +0.038] & $\mathbf{+0.091}$ [+0.022, +0.164] \\
    ViT-g single $-$ SalFoM           & $\mathbf{+0.036}$ [+0.026, +0.048] & $\mathbf{+0.130}$ [+0.060, +0.205] \\
    ViT-g ensemble $-$ SalFoM         & $\mathbf{+0.039}$ [+0.029, +0.050] & $\mathbf{+0.160}$ [+0.089, +0.233] \\
    ViT-g single $-$ ViT-L single     & $\mathbf{+0.017}$ [+0.012, +0.022] & $\mathbf{+0.108}$ [+0.075, +0.142] \\
    ViT-g ensemble $-$ ViT-L ensemble & $\mathbf{+0.012}$ [+0.008, +0.016] & $\mathbf{+0.069}$ [+0.041, +0.097] \\
    \bottomrule
  \end{tabular}
\end{table}

\begin{figure}[htbp]
  \centering
  \includegraphics[width=\textwidth]{val_delta.pdf}
  \caption{Per-video CC difference between the SURGE-D ViT-g ensemble and SalFoM
  on the 100 DHF1K validation videos, sorted. The dashed line is the mean
  difference; the interval is the paired bootstrap 95\,\% interval}\label{fig:valdelta}
\end{figure}

Two qualifications apply. First, SURGE-D uses frames on both sides of the
predicted frame (\cref{sec:formulation}), whereas SalFoM predicts the last frame
of its input clip; part of the difference may therefore come from access to
future context, which SalFoM does not use. Second, the comparison involves one
re-scored baseline. The published validation scores in the lower block of
\cref{tab:dhf1k} were computed with other code, which agrees with the present
evaluator to within 0.0003 CC on the reference maps
(\cref{sec:metric-validation}); with that caveat, SURGE-D's validation CC is
0.046--0.049 above the highest of them. Reported validation scores also differ
between papers: UniST~\cite{xiong2025unist}, for instance, lists TMFI-Net at
0.524 rather than 0.554 and reports 0.541 for itself. Several other
peer-reviewed models published since 2024~\cite{li2025tm2sp,zhang2024msff,moradi2026lstd}
report DHF1K results whose tables could not be accessed or that refer to the
test set; test-set scores cannot be compared with the validation scores of
this thesis.

\paragraph{DHF1K test set.} Maps for the 300 official test videos were produced
with the ViT-g three-seed ensemble and checked for format by passing five
validation videos through the same export, which reproduced their validation
scores to within $10^{-4}$ CC. The scores on the test set are computed by the
dataset authors on request and were not available when this thesis was
written, so no test-set result is claimed.

\subsection{Generalisation to Other Datasets}\label{sec:transfer}

\paragraph{UCF Sports.} UCF Sports differs from DHF1K in two ways: its gaze was
recorded during an action-recognition task~\cite{mathe2015actions}, and it is
small (103 training videos). Without any UCF training, the DHF1K model reached
a CC of 0.608 (\cref{tab:ucf}), well below the fine-tuned models, which confirms
that the task-driven gaze of this dataset differs from free viewing. After
fine-tuning, the ViT-g ensemble reached CC 0.735 and the best NSS, SIM and AUC-J
of the SURGE-D models.

\begin{table}[htbp]
  \centering
  \caption{Results on the UCF Sports test set (47 videos, 2,929 frames). Published
  results were computed with other evaluation code}\label{tab:ucf}
  \small
  \begin{tabular}{@{}lcccccc@{}}
    \toprule
    \textbf{Model} & \textbf{CC}$\uparrow$ & \textbf{NSS}$\uparrow$ & \textbf{SIM}$\uparrow$
      & \textbf{KLD}$\downarrow$ & \textbf{AUC-J}$\uparrow$ & \textbf{sAUC}$\uparrow$ \\
    \midrule
    SURGE-D ViT-L, zero-shot                & 0.6077 & 3.128 & 0.5033 & 1.113 & 0.9223 & 0.7453 \\
    SURGE-D ViT-L, fine-tuned               & 0.7252 & 3.916 & 0.5799 & 0.773 & 0.9406 & 0.8089 \\
    SURGE-D ViT-L, 3-seed ensemble          & 0.7331 & 4.028 & 0.5833 & 0.749 & 0.9430 & 0.8153 \\
    SURGE-D ViT-g, fine-tuned               & 0.7305 & 4.061 & 0.5831 & 0.747 & 0.9442 & 0.8145 \\
    SURGE-D ViT-g, 3-seed ensemble          & 0.7350 & 4.085 & 0.5857 & \textbf{0.736} & \textbf{0.9449} & \textbf{0.8173} \\
    \midrule
    \multicolumn{7}{@{}l}{\emph{Published (\cite{girmaji2025minimalistic}; SalFoM from~\cite{moradi2024salfom})}} \\
    ViNet-E~\cite{girmaji2025minimalistic}   & \textbf{0.744} & \textbf{4.156} & \textbf{0.587} & -- & 0.941 & -- \\
    ViNet-A~\cite{girmaji2025minimalistic}   & 0.734 & 4.108 & 0.586 & -- & 0.940 & -- \\
    THTD-Net~\cite{moradi2024transformer}    & 0.711 & 3.840 & 0.565 & -- & 0.933 & -- \\
    TMFI-Net~\cite{zhou2023transformer}      & 0.707 & 3.863 & 0.565 & -- & 0.936 & -- \\
    SalFoM~\cite{moradi2024salfom}           & 0.631 & 3.543 & 0.516 & -- & 0.928 & -- \\
    \bottomrule
  \end{tabular}
\end{table}

On UCF Sports the encoder made a smaller difference than on DHF1K: the ViT-g
ensemble was not significantly better than the ViT-L ensemble
($\Delta$CC $+0.0019$ $[-0.0071, +0.0125]$), whereas the three-seed ViT-L
ensemble was significantly better than its single member ($+0.0080$
$[+0.0039, +0.0123]$). Against published results, SURGE-D clearly exceeds
SalFoM (by about 0.10 CC) and TMFI-Net, ties ViNet-A ($+0.001$) and remains
below ViNet-E ($-0.009$, a tie under the rule of \cref{tab:rules}, but below
ViNet-E on CC, NSS and SIM, though above it on AUC-J). SURGE-D is therefore competitive on UCF Sports but not ahead of the
best published model. Both ViNet variants use encoders pretrained for action
recognition or action localisation~\cite{girmaji2025minimalistic}, which may
suit the task-driven gaze of this dataset.

\paragraph{Audio-visual datasets.} \Cref{tab:av} reports the six datasets of the
STAViS benchmark~\cite{tsiami2020stavis}, as fold means. The DHF1K model already
transferred reasonably without any training (zero-shot row), and fine-tuning
improved every dataset, most of all Coutrot2 ($+0.17$ CC), whose meetings
differ most from DHF1K. With ViT-g, the ensemble was significantly better than
the ViT-L ensemble on DIEM, Coutrot1, ETMD and AVAD (paired bootstrap over the
pooled test videos), not significantly different on Coutrot2, and tied on SumMe.

\begin{table}[htbp]
  \centering
  \caption{CC on the six audio-visual datasets (fold means; DIEM has one split).
  The lower block lists the best peer-reviewed results whose tables could be
  verified, computed with other evaluation code. Visual-only: ViNet-E, or ViNet-A
  on Coutrot2~\cite{girmaji2025minimalistic}; on SumMe, which those models do not
  report, the visual-only variant of the relevance-guided fusion model AVRSP of
  Yu et al.~\cite{yu2026relevance}. Audio-visual: AVRSP~\cite{yu2026relevance} on
  DIEM and Coutrot2, DiffSal~\cite{xiong2024diffsal} on the other four}\label{tab:av}
  \small
  \begin{tabular}{@{}lcccccc@{}}
    \toprule
    \textbf{Model} & \textbf{DIEM} & \textbf{Coutrot1} & \textbf{Coutrot2} & \textbf{ETMD}
      & \textbf{AVAD} & \textbf{SumMe} \\
    \midrule
    SURGE-D ViT-L, zero-shot       & 0.675 & 0.600 & 0.645 & 0.588 & 0.692 & 0.551 \\
    SURGE-D ViT-L, fine-tuned      & 0.706 & 0.632 & 0.816 & 0.662 & 0.718 & 0.581 \\
    SURGE-D ViT-L, 3-seed ensemble & 0.714 & 0.641 & 0.807 & 0.670 & 0.729 & \textbf{0.593} \\
    SURGE-D ViT-g, fine-tuned      & 0.714 & 0.648 & 0.820 & 0.672 & 0.731 & 0.586 \\
    SURGE-D ViT-g, 3-seed ensemble & \textbf{0.729} & \textbf{0.657} & 0.819 & \textbf{0.676} & \textbf{0.740} & 0.590 \\
    \midrule
    Best visual-only (published)   & 0.701 & 0.614 & \textbf{0.862} & 0.632 & 0.729 & 0.494 \\
    Best audio-visual (published)  & 0.685 & 0.638 & 0.837 & 0.652 & 0.738 & 0.572 \\
    \midrule
    ViT-g ens.\ $-$ best visual-only  & $+0.028$ & $+0.043$ & $-0.043$ & $+0.044$ & $+0.011$ & $+0.096$ \\
    ViT-g ens.\ $-$ best audio-visual & $+0.044$ & $+0.019$ & $-0.018$ & $+0.024$ & $+0.002$ & $+0.018$ \\
    \bottomrule
  \end{tabular}
\end{table}

Against the best peer-reviewed results that could be verified, the ViT-g
ensemble has the higher CC on five of the six datasets. The margin is clear on
DIEM, Coutrot1, ETMD and SumMe. On AVAD, the margin over the best visual-only
model is at the threshold ($+0.011$), and the result ties the best audio-visual
model. On Coutrot2 it is lower than both. The other metrics broadly agree with
CC: the ensemble's NSS (DIEM 2.88, Coutrot1 3.30, ETMD 3.69, AVAD 4.31, SumMe
3.26) is above the published values of both reference models on these five
datasets, while SIM is lower than DiffSal's on Coutrot1 and AVAD.

These comparisons have three limits. They cross evaluation code. They are
restricted to peer-reviewed results, so results that had not appeared in
peer-reviewed venues when this thesis was written are not included. And they
are restricted to results whose published tables could be accessed: further
recent audio-visual models, among them TM2SP~\cite{li2025tm2sp},
FOSal~\cite{yu2026fosal} and the model of Mao et al.~\cite{mao2026audio}, could
not be checked and are therefore not represented. The table should consequently
not be read as establishing a state of the art on these datasets. Coutrot2 is
discussed in \cref{sec:discussion}.

\paragraph{Hollywood-2 with the Vig gaze data.} \Cref{tab:hollywood} reports the
884 test clips of Hollywood-2 with the gaze of Vig et al.~\cite{vig2012space}
(257,733 scored frames). These numbers are not comparable with published
Hollywood-2 results (\cref{sec:dataset}). With this many test clips, small
differences are resolved precisely: ViT-g was better than ViT-L on all six
metrics as a single model ($\Delta$CC $+0.0046$ $[+0.0032, +0.0060]$) and on
five as an ensemble ($+0.0041$ $[+0.0030, +0.0052]$). Ensembling raised CC, KLD,
AUC-J and sAUC but significantly \emph{lowered} NSS and SIM, for both encoders.
With at most five observers per frame, the ground truth is sparse, and the
averaging of three members smooths the peaks that NSS and SIM reward.

\begin{table}[htbp]
  \centering
  \caption{Results on Hollywood-2 with the gaze of Vig et al.\ (884 test clips).
  Not comparable with published Hollywood-2 results, which use other ground
  truth}\label{tab:hollywood}
  \small
  \begin{tabular}{@{}lcccccc@{}}
    \toprule
    \textbf{Model} & \textbf{CC}$\uparrow$ & \textbf{NSS}$\uparrow$ & \textbf{SIM}$\uparrow$
      & \textbf{KLD}$\downarrow$ & \textbf{AUC-J}$\uparrow$ & \textbf{sAUC}$\uparrow$ \\
    \midrule
    SURGE-D ViT-L, zero-shot        & 0.632 & 3.595 & 0.503 & 1.061 & 0.951 & 0.780 \\
    SURGE-D ViT-L, fine-tuned       & 0.6798 & 3.898 & 0.5178 & 0.925 & 0.9564 & 0.8167 \\
    SURGE-D ViT-L, 3-seed ensemble  & 0.6841 & 3.875 & 0.5073 & 0.918 & 0.9574 & 0.8192 \\
    SURGE-D ViT-g, fine-tuned       & 0.6844 & \textbf{3.933} & \textbf{0.5214} & 0.917 & 0.9569 & 0.8180 \\
    SURGE-D ViT-g, 3-seed ensemble  & \textbf{0.6882} & 3.907 & 0.5104 & \textbf{0.911} & \textbf{0.9577} & \textbf{0.8199} \\
    \bottomrule
  \end{tabular}
\end{table}

\subsection{Qualitative Results}\label{sec:qualitative}

\Cref{fig:qualitative} shows frames from five DHF1K validation videos, chosen by
a rule fixed before any map was inspected: the videos at the 10th, 50th and
90th percentile of the ensemble's per-video CC, and the videos with the largest
gain and the largest loss relative to SalFoM; the frame shown is the middle
frame of each video.

\begin{figure}[htbp]
  \centering
  \includegraphics[width=\textwidth]{qualitative.pdf}
  \caption{Qualitative examples on DHF1K validation. Rows: videos at the 10th,
  50th and 90th percentile of SURGE-D's per-video CC, and the videos with the
  largest gain (0696) and loss (0663) relative to SalFoM; columns: middle frame
  with the observers' fixations, ground-truth density, SalFoM and the SURGE-D
  ViT-g ensemble, overlaid on the frame. Numbers give the CC of the shown frame.
  Video-level CC (SURGE-D / SalFoM): 0650 0.455/0.376, 0632 0.613/0.639, 0685
  0.741/0.715, 0696 0.664/0.437, 0663 0.388/0.462}\label{fig:qualitative}
\end{figure}

On these frames, the two models produce similar maps: both locate the horse
rider, the surfer and the flying object that observers follow, and on the
single frames shown SalFoM is in fact marginally closer to the ground truth in
four of the five rows and equal in the fifth. Single frames are a poor guide to
average differences that amount to a few hundredths of CC: for video 0696, for
example, SURGE-D's CC over the whole video is 0.23 higher, yet the two maps at
its middle frame are almost equal. The last row shows a characteristic failure.
In a panel discussion, the observers' fixations concentrate on one person on
the right, while SURGE-D spreads its prediction over several faces at the
table. In conversation scenes, gaze is drawn to the faces of the people who are
speaking~\cite{coutrot2014how}, which a model that sees only the image cannot
determine; the same limitation reappears on Coutrot2 below.

\subsection{Computational Cost}\label{sec:cost}

\Cref{tab:cost} reports the cost of SURGE-D, measured with the PyTorch FLOP
counter on one $16\times256\times256$ clip and timed on one DHF1K validation
video with frames preloaded into memory (bfloat16, network forward pass only).
Because one forward pass yields maps for 16 original frames, the cost per output
frame is one sixteenth of the cost per clip.

\begin{table}[htbp]
  \centering
  \caption{Computational cost of SURGE-D on an RTX A5000 (forward pass only).
  TFLOPs count multiply--add pairs as two operations}\label{tab:cost}
  \small
  \begin{tabular}{@{}lcc@{}}
    \toprule
    & \textbf{ViT-L} & \textbf{ViT-g} \\
    \midrule
    Encoder parameters (frozen)                & 303.9\,M & 1,012.2\,M \\
    Decoder parameters (trained)               & 6.99\,M & 7.39\,M \\
    Encoder TFLOPs per clip                    & 1.656 & 5.088 \\
    Decoder TFLOPs per clip                    & 0.166 & 0.168 \\
    TFLOPs per output frame                    & 0.114 & 0.328 \\
    Throughput, batch 4 (output frames/s)      & 301 & 129 \\
    Latency per clip, batch 1 (ms)             & 116 & 163 \\
    Peak GPU memory, batch 4 (GB)              & 2.8 & 7.1 \\
    \bottomrule
  \end{tabular}
\end{table}

Both models run faster than real time on this GPU. For comparison, the UMT-L
encoder that SalFoM fine-tunes requires 1.90 TFLOPs for one $16\times224\times224$
clip (3,136 tokens; counted on a re-implementation loaded from the released
checkpoint), and SalFoM runs it once per output frame in a sliding
window~\cite{moradi2024salfom}. Counting the encoder alone, SalFoM therefore needs
about 17 times the computation per output frame of SURGE-D with ViT-L, and
about 6 times that of SURGE-D with ViT-g. The multi-frame output thus allows
an encoder of three times the size at a fraction of the cost of a per-frame
design.

%----------------------------------------------------------------------------
\section{Ablation Study}\label{sec:ablation}

\Cref{tab:ablation} collects the ablations. Each was evaluated on the holdout
under the rules of \cref{tab:rules}; for single-seed runs, the reference is the
frozen ViT-L model of seed 0 trained with the same three-epoch schedule.

\begin{table}[htbp]
  \centering
  \caption{Ablations on the DHF1K holdout (60 videos) unless marked
  \enquote{val}. $\Delta$CC is relative to the reference in each block}\label{tab:ablation}
  \small
  \begin{tabularx}{\textwidth}{@{}>{\raggedright}X ccl@{}}
    \toprule
    \textbf{Variant} & \textbf{CC} & $\Delta$\textbf{CC} & \textbf{Outcome under the rule} \\
    \midrule
    \multicolumn{4}{@{}l}{\emph{Components (ViT-L, seed 0; reference CC 0.5772, NSS 3.329)}} \\
    Without centre prior                 & 0.5772 & $-0.0000$ & no difference \\
    Without auxiliary heads              & 0.5757 & $-0.0015$ & no difference \\
    With SALICON co-training             & 0.5801 & $+0.0028$ & no difference \\
    Stride 4 ($\approx$7.5\,fps, 2\,s window) & 0.5803 & $+0.0031$ & difference (at the threshold) \\
    \midrule
    \multicolumn{4}{@{}l}{\emph{Encoder adaptation (ViT-L, three seeds)}} \\
    LoRA, seed 0 / 1 / 2 & 0.5743 / 0.5791 / 0.5739 & $-0.0029$ / $-0.0010$ / $-0.0037$ & rejected (lost 3 of 3) \\
    \midrule
    \multicolumn{4}{@{}l}{\emph{Encoder scale (seed 0; gate 0.5833)}} \\
    ViT-g instead of ViT-L                 & 0.5866 & $+0.0083$ vs ViT-L mean & adopted \\
    \midrule
    \multicolumn{4}{@{}l}{\emph{Clip length (ViT-g, seed 0; gate 0.5904)}} \\
    32 frames instead of 16                & 0.5902 & $+0.0036$ vs seed 0 & not pursued (gate missed by 0.0002) \\
    \midrule
    \multicolumn{4}{@{}l}{\emph{Training data (ViT-g, val, three-seed ensembles)}} \\
    600 instead of 540 training videos (val) & 0.6044 & $+0.0015$ [$-$0.0006, +0.0037] & not adopted \\
    \bottomrule
  \end{tabularx}
\end{table}

\paragraph{Encoder adaptation.} LoRA was expected to help, because adapting a
pretrained encoder to the target task usually improves on frozen features, and
LoRA does so while training only a small fraction of the parameters that full
fine-tuning would change~\cite{hu2022lora}. Here it did not. Each LoRA model started from the frozen model of the same seed and lost to
it on the holdout in all three seeds (mean $-0.0025$ CC), and its best epoch was
always the first. With only 540 training videos, adapting 1.6 million encoder
parameters appears to add more variance than signal; the same early
overfitting that limits the decoder (\cref{fig:g1}) would act on the adapted
encoder as well. Full fine-tuning, as done by SalFoM, was not tested.

\paragraph{Components.} Removing the learned centre prior or the auxiliary
heads, and adding 10,000 static SALICON images per epoch, changed holdout CC by
less than the threshold of 0.003. The centre prior's lack of effect suggests
that the encoder features and the convolutional decoder can already represent
the centre bias, for instance through the image borders, which the squashed
input keeps at fixed positions; the explicit prior then adds nothing that can
be measured. The static images,
which helped CNN-era models~\cite{wang2018revisiting}, appear redundant once
the encoder has been pretrained on large amounts of video. A stride of four,
which halves the cost per output frame, reached $+0.0031$, exactly at the
threshold; with one seed, this is a candidate for future work rather than a
finding.

\paragraph{Encoder scale.} Replacing ViT-L by ViT-g, with every other setting
unchanged, raised holdout CC for all three seeds (0.5866, 0.5867 and 0.5829
against 0.5772, 0.5801 and 0.5776) and validation CC by 0.017 for single models
(\cref{tab:dhf1k-boot}). Scale was the largest single improvement found after
the choice of encoder family.

\paragraph{Temporal context and training data.} Doubling the clip to 32 frames
(about 2.1\,s) reached a holdout CC of 0.5902, 0.0002 below its gate, so it was
not pursued; unlike the 16-frame runs, it was still improving at its last
epoch, so a longer schedule might change this conclusion, but that would be a
new experiment. Training on all 600 DHF1K training videos instead of 540 gave
validation scores within noise of the 540-video models. Neither more context
nor 11\,\% more videos made a measurable difference.

\paragraph{Ensembles and output blur.} Averaging three seeds raised validation
CC by 0.007 for ViT-L and by 0.003 for ViT-g (\cref{tab:dhf1k}); the stronger
encoder leaves less variance for the ensemble to remove. \Cref{fig:sigma} shows
that the optimal blur is broad: holdout CC varied by at most 0.003 between
$\sigma=0$ and $\sigma=12$ for every configuration. The ViT-L ensemble preferred
less blur than its single model (8 against 12 pixels), consistent with the
smoothing that averaging already provides.

\begin{figure}[htbp]
  \centering
  \includegraphics[width=0.62\textwidth]{sigma.pdf}
  \caption{Holdout CC as a function of the Gaussian blur applied to the output,
  for single models and three-seed ensembles. The selected values are 12 (ViT-L
  single) and 8 (others)}\label{fig:sigma}
\end{figure}

\paragraph{Audio.}\label{sec:audio} The audio fusion module of
\cref{sec:audio-fusion} was trained on every fold of the six audio-visual
datasets with the same recipe as the visual fine-tuning (ViT-L, seed 0).
\Cref{tab:audio} shows that it did not improve the visual model on any dataset:
the largest change was $+0.001$ CC, far below the threshold of $+0.010$. The
control is informative nonetheless. On Coutrot2, scoring the trained model with
the wrong soundtrack cost 0.018 CC, so the model had learned to use the audio
there, but with the correct soundtrack it only matched the visual model, which
already finds the faces in these meeting scenes from the image alone. With five
test videos per fold, this observation is suggestive rather than conclusive.

\begin{table}[htbp]
  \centering
  \caption{Audio fusion on the audio-visual datasets (ViT-L, fold-mean CC). The
  rule requires audio $-$ visual $\geq+0.010$ and matched $-$ rotated
  $\geq+0.005$}\label{tab:audio}
  \small
  \begin{tabular}{@{}lccccc@{}}
    \toprule
    \textbf{Dataset} & \textbf{Visual} & \textbf{Audio, matched} & \textbf{Audio, rotated}
      & \textbf{Audio $-$ visual} & \textbf{Matched $-$ rotated} \\
    \midrule
    DIEM     & 0.706 & 0.706 & 0.706 & $+0.0008$ & $+0.0002$ \\
    Coutrot1 & 0.632 & 0.633 & 0.630 & $+0.0010$ & $+0.0031$ \\
    Coutrot2 & 0.816 & 0.816 & 0.797 & $-0.0003$ & $+0.0183$ \\
    ETMD     & 0.662 & 0.661 & 0.661 & $-0.0014$ & $-0.0001$ \\
    AVAD     & 0.718 & 0.718 & 0.717 & $+0.0002$ & $+0.0013$ \\
    SumMe    & 0.581 & 0.579 & 0.580 & $-0.0021$ & $-0.0011$ \\
    \bottomrule
  \end{tabular}
\end{table}

%----------------------------------------------------------------------------
\section{Discussion}\label{sec:discussion}

This section interprets the results in terms of the research questions of
\cref{sec:questions}.

\paragraph{RQ1: Can a frozen self-supervised video encoder with a light decoder
predict fixations better than the strongest foundation-model approach?} On
DHF1K, yes, within the limits of the available test data. With ViT-L, a frozen
V-JEPA~2 encoder and a 7-million-parameter decoder already exceeded SalFoM, a
model whose 300-million-parameter encoder is fully fine-tuned, under an
identical evaluation, on five of six metrics as a single model and on all six as
an ensemble; with ViT-g, the margin of the single model grew to $+0.036$ CC and
covered all six metrics. Two features of the result matter for
its interpretation. First, the gain does not come from adapting the encoder:
the frozen encoder outperformed its LoRA-adapted version, and the result
therefore speaks to the quality of the pretrained V-JEPA representations
rather than to a training technique. The gate G1 comparison, in which
V-JEPA~2 outperformed UMT-L by 0.032 CC under the same frozen protocol,
supports this reading, since UMT-L is the encoder family that SalFoM builds on.
Second, the evidence is a validation-set comparison with one re-scored baseline
and SURGE-D's access to future frames; the official test-set evaluation is
still outstanding.

\paragraph{RQ2: Does predictive surprise carry information about where people
look?} Not in the form available from V-JEPA~2. Its predictor's error was
anti-correlated with fixations at 15\,fps and was beaten by both trivial
reference cues. The diagnostics point to the reason: the error tracks the
atypicality of each region rather than violated expectations, and the
pre-training masks give the predictor no reason to forecast. This is a
statement about one released predictor and one way of reading it out, not a
refutation of surprise-based accounts of attention~\cite{itti2009bayesian};
a predictor trained to extrapolate in time could behave differently.

\paragraph{RQ3: Which design factors matter when the encoder is frozen?} The
results separate the factors cleanly. The encoder dominates: its pre-training
(V-JEPA~2 over V-JEPA~2.1 and UMT-L) and its scale (ViT-g over ViT-L) produced
the only gains that were large relative to seed noise on DHF1K. Seed ensembles
added a smaller but reliable gain. Nearly everything else, including LoRA, the
learned centre prior, deep supervision, static-image co-training, 32-frame
clips and 11\,\% more training videos, made no measurable difference. A
consistent explanation is that, with a strong frozen encoder, the limiting
factor is the small number of training videos: the decoder overfits within
three epochs (\cref{fig:g1}), and the additions that would increase its
effective capacity do not help, whereas better features do.

\paragraph{RQ4: How well does the approach transfer, and does audio help?} The
model transferred well. Fine-tuned for three epochs, it matched or exceeded the
best peer-reviewed visual-only and audio-visual results that could be verified
on five of the six audio-visual datasets, and it was competitive on UCF Sports,
where action-oriented encoders retain an edge. The one clear failure, Coutrot2,
has a simple account: in recordings of conversations, observers look at the
person who is speaking, which an image-only model cannot determine, as the last
row of \cref{fig:qualitative} also illustrates. The audio experiment shows that
adding the soundtrack in a simple way does not resolve this. The fusion module
learned to depend on the audio on Coutrot2, but it did not improve on the visual
model. This agrees with the observations of Jain et al.~\cite{jain2021vinet} and
Agrawal et al.~\cite{agrawal2022does}, who found that deep audio-visual saliency
models tend to become insensitive to their audio input and that none of a range
of audio modules and fusion techniques improved on visual-only models across
ten datasets. Exploiting sound in this setting probably requires explicit
speaker localisation rather than global feature modulation.

\paragraph{Relation to the original proposal.} The final model differs from the
proposed one in three respects, each decided by a pre-registered test: the
encoder is V-JEPA~2 rather than V-JEPA~2.1, it is frozen rather than adapted with
LoRA, and the surprise branch that gave the project its name is absent. The
proposal's ranking of what drives performance, which placed the encoder first,
was confirmed. Its expectation that LoRA would help, based on results with much
larger training sets, was not.

%----------------------------------------------------------------------------
\section{Threats to Validity}\label{sec:threats}

\paragraph{Internal validity.} All selection used holdout data, every test set
was scored once, and every rule was written in advance, which protects against
selection bias. Two residual weaknesses remain. The component ablations used
one seed, and they were judged against a fixed threshold rather than a
statistical test, so effects of a few thousandths of CC cannot be excluded.
In addition, several thresholds (for example the margin of 0.005 in the gates)
were chosen by judgement from the observed seed spread rather than derived
formally.

\paragraph{Construct validity.} The six metrics capture different aspects of a
saliency map~\cite{bylinskii2019metrics} and did not always agree; ensembles,
for instance, improved CC but lowered NSS and SIM on the sparse Hollywood-2
(Vig gaze) data. The use of one map for all metrics, with the blur chosen by CC,
favours CC slightly over the other metrics.

\paragraph{External validity and comparability.} Comparisons with published
numbers cross evaluation code and, in places, protocol details, such as the
handling of frames without fixations or the exact fold composition. Only one
baseline, SalFoM, was re-scored under the protocol of this thesis. SURGE-D's
predictions use about 0.27\,s or more of future context, which causal models do
not have. The DHF1K test-set result and the standard Hollywood-2 ground truth
were not available. Finally, the pre-training data of V-JEPA~2 include large
collections of internet video, and DHF1K is drawn from YouTube; an overlap
between pre-training and test videos cannot be ruled out. The same is true for
every method that relies on a large pretrained video encoder, including SalFoM.
